\chapter{Fazit und Ausblick}
\label{fazit}

Das Ziel dieser Studienarbeit war die Entwicklung und das Testen grundlegender Fahrassistenzsysteme auf einem ferngesteuerten Modellfahrzeug. Im Rahmen des Projekts wurden erfolgreich eine Notbremsfunktion, eine Linienverfolgung sowie eine Fernsteuerung über eine Weboberfläche implementiert. Die Arbeit liefert dabei Antworten auf zentrale Forschungsfragen bezüglich der erforderlichen Hardware, der Zuverlässigkeit von Sensoren und der Übertragbarkeit der Ergebnisse. 

Die Schildererkennung konnte aufgrund von Zeitmangel und der Komplexität der Implementierung nicht realisiert werden, weshalb diese in Zukunft implementiert werden könnte.

Die erste Forschungsfrage befasste sich damit, welche Hardwarekomponenten grundlegend erforderlich sind. Dabei hat die Implementierung hat gezeigt, dass eine minimale, aber effektive Konfiguration aus einer zentralen Steuereinheit, wie einem Raspberry Pi, Motoren mit Motortreibern und einer Reihe von Sensoren besteht. Unverzichtbar waren dabei Ultraschallsensoren für die Hinderniserkennung, ein Hall-Sensor zur Geschwindigkeitsmessung und ein Kameramodul für die Linienverfolgung. Eine stabile Stromversorgung wie in diesem Fall mittels \ac{LiPo}-Akku und Spannungsreglern ist ebenfalls essenziell.

Die zweite Frage untersuchte die Zuverlässigkeit von Ultraschallsensoren. Die durchgeführten Tests ergaben aufschlussreiche Ergebnisse: Bei harten, glatten Oberflächen wie Wänden oder Holz lieferten die Sensoren über verschiedene Distanzen von 50 cm bis 200 cm sehr präzise und zuverlässige Messwerte. Bei schallabsorbierenden Materialien wie Schaumstoff versagte die Messung hingegen vollständig. Objekte mit unregelmäßigen Oberflächen, wie ein Eierkarton, führten teilweise zu einer Streuung der Messwerte und machten eine genaue Abstandsmessung schwierig. \\
Die Zuverlässigkeit ist also stark vom Material des Hindernisses abhängig. Außerdem ist der Winkel der Sensoren zum Hindernis ein kritischer Faktor, da schräg auftreffende Ultraschallwellen nur teilweise oder gar nicht zum Sensor zurückreflektiert werden, was zu ungenauen Messungen führen kann. Insgesamt sind Ultraschallsensoren für die Hinderniserkennung unter kontrollierten Bedingungen durchaus zuverlässig, ihre Effektivität kann jedoch durch Material des Hindernisses und die Anordnung der Sensoren stark beeinträchtigt werden. 

Die dritte Forschungsfrage zielte auf die Faktoren ab, die die Genauigkeit der Systeme beeinflussen. Wie bereits bei den Ultraschallsensoren erwähnt ist das Material von Hindernissen ein wichtiger Faktor, da die Notbremsfunktion auf die Werte der Ultraschallsensoren angewiesen ist. Sind diese nicht korrekt, kann es zu einer Kollision kommen. Auch das Material der Reifen und des Untergrund haben Auswirkungen auf die Effektivität der Notbremsfunktion, als auch der Manövrierfähigkeit. Die Genauigkeit der Linienverfolgung wurde zudem von den Lichtverhältnissen und der Linienbreite beeinflusst. Eine Optimierung kann durch die Kombination verschiedener Sensortypen erreicht werden, um die Schwächen eines Sensortyps auszugleichen. 

Die letzte Frage behandelte die Übertragbarkeit der Ergebnisse auf reale Fahrzeuge. Bei den in dieser Arbeit entwickelten Konzepte und Algorithmen ist eine Skalierung prinzipiell möglich. Die grundlegende Logik einer Notbremsung basierend auf Abstand und Geschwindigkeit oder die Bildverarbeitung zur Linienverfolgung sind in realen Fahrzeugen ähnlich. Die Herausforderungen bei der Übertragung liegen jedoch in der Komplexität und den Sicherheitsanforderungen der realen Welt. Echte Fahrzeuge erfordern eine deutlich robustere Sensorik, weshalb bei der Verwendung von Ultraschallsensoren mehrere notwendig wären, um zuverlässig arbeiten zu können. Generell sind redundante Systeme zur Ausfallsicherheit wichtig, um eine gewisse Sicherheit zu gewährleisten. In einem realen System müssen die Komponenten außerdem mit einer größeren Vielfalt an unvorhersehbaren Verkehrssituationen, Wetterbedingungen und Geschwindigkeiten umgehen können. Deshalb wären zusätzliche Tests und Validierungen erforderlich, um die Funktionalität bei Regen und anderen Wetterbedingungen zu gewährleisten. Auch die Linienverfolgung müsste für höhere Geschwindigkeiten und komplexere Straßenverhältnisse optimiert werden.

Die durch das Projekt gewonnenen Erkenntnisse bilden eine wertvolle Grundlage, können jedoch nicht alle Bedingungen und Situationen, die in einem realen System auftreten, abdecken. Das Testen von Fahrassistenzsystemen in einem kleineren Maßstab ist daher nur bedingt möglich. Um die Übertragbarkeit zu verbessern, müsste beispielsweise der Maßstab des Modells erhöht und die Anzahl und Vielfalt der Sensoren erweitert werden, um eine realistischere Simulation zu ermöglichen.

% \section*{Leitfragen für das Fazit}
% \begin{itemize}
%     \item Welche zentralen Ergebnisse wurden in der Arbeit erzielt?
%     \item Inwiefern wurden die Forschungsfragen beantwortet?
%     \item Welche Herausforderungen traten während der Arbeit auf und wie wurden sie gelöst?
%     \item Welche Bedeutung haben die Ergebnisse für das behandelte Thema?
%     \item Welche Grenzen und Einschränkungen hat die Arbeit?
% \end{itemize}

% \section*{Leitfragen für den Ausblick}
% \begin{itemize}
%     \item Welche offenen Fragen bleiben nach der Arbeit bestehen?
%     \item Welche weiterführenden Forschungsansätze ergeben sich aus den Ergebnissen?
%     \item Wie könnten die Ergebnisse in der Praxis angewendet werden?
%     \item Welche zukünftigen Entwicklungen könnten das Thema beeinflussen?
%     \item Welche Verbesserungen oder Erweiterungen wären für die Arbeit denkbar?
% \end{itemize}
